\documentclass[11pt]{article}
\usepackage[margin=1in]{geometry}
\usepackage{amsmath,amssymb,amsthm}
\usepackage{graphicx}
\usepackage[hidelinks]{hyperref}

\newtheorem{theorem}{Theorem}
\newtheorem{lemma}{Lemma}

\newcommand{\eps}{\varepsilon}
\newcommand{\norm}[1]{\left\lVert #1 \right\rVert}
\newcommand{\R}{\mathbb{R}}
\newcommand{\Sph}{S_X}

\title{A Counterexample to the Milman-Moduli Quasi-Orthogonality Conjecture}
\author{Run fa\_banach\_001, agent\_lane\_10}
\date{June 16, 2026}

\begin{document}
\maketitle

\section*{Source question}

The source is Grigory Ivanov and Horst Martini, \emph{New Moduli for Banach
Spaces}, arXiv:1609.01587.  In Section 6, after defining Milman's moduli
\[
\beta^-_X(\eps)=\inf_{x,y\in S_X}
\left(\max\{\norm{x+\eps y},\norm{x-\eps y}\}-1\right)
\]
and
\[
\beta^+_X(\eps)=\sup_{x,y\in S_X}
\left(\min\{\norm{x+\eps y},\norm{x-\eps y}\}-1\right),
\]
the authors conjecture that it is enough to restrict to pairs with
\(y\) quasi-orthogonal to \(x\).  With their notation
\[
\zeta^-_X(\eps)=\inf\{\norm{x+\eps y}:x,y\in S_X,\ y\urcorner x\},
\quad
\zeta^+_X(\eps)=\sup\{\norm{x+\eps y}:x,y\in S_X,\ y\urcorner x\},
\]
the conjectured identities are
\[
\zeta^-_X(\eps)-1=\beta^-_X(\eps),\qquad
\zeta^+_X(\eps)-1=\beta^+_X(\eps)
\]
for positive \(\eps\).  Here \(y\urcorner x\) means that some
\(p\in J_1(x)\) satisfies \(p(y)=0\).

\begin{center}
\includegraphics[width=0.94\textwidth]{figures/open_problem_crop.png}
\end{center}

\section*{Result}

\begin{theorem}
The first identity in the conjecture is false.  In the real Banach space
\(X=\ell_4^2\),
\[
\zeta^-_X(1/2)-1<\beta^-_X(1/2).
\]
\end{theorem}

\section*{Idea of the proof}

The quantity \(\beta^-_X\) is protected by a two-sign maximum: for every pair
\((u,v)\), at least one of \(\norm{u+\eps v}\) and \(\norm{u-\eps v}\) must be
large.  The quantity \(\zeta^-_X\) only asks about the single chosen sign
\(\norm{x+\eps y}\), even when \(y\) lies in a supporting hyperplane at \(x\).
In \(\ell_4^2\), a point \(x\) very close to a coordinate vector has a tangent
direction whose negative second coordinate moves slightly inward for the
positive sign \(x+\eps y\).  This lowers the one-sign value below the
two-sign Milman lower bound.

\section*{Proof}

Work in \(X=(\R^2,\norm{\cdot}_4)\), where
\[
\norm{(a,b)}_4=(|a|^4+|b|^4)^{1/4}.
\]

\begin{lemma}
For every \(\eps>0\),
\[
\beta^-_{\ell_4^2}(\eps)+1=(1+\eps^4)^{1/4}.
\]
\end{lemma}

\begin{proof}
For real numbers \(a,b\),
\[
|a+b|^4+|a-b|^4=2(a^4+6a^2b^2+b^4)\geq 2(a^4+b^4).
\]
Hence, if \(u,v\in S_{\ell_4^2}\), then
\[
\norm{u+\eps v}_4^4+\norm{u-\eps v}_4^4
=2(1+\eps^4)+12\eps^2\sum_i u_i^2v_i^2
\geq 2(1+\eps^4).
\]
Therefore
\[
\max\{\norm{u+\eps v}_4,\norm{u-\eps v}_4\}\geq (1+\eps^4)^{1/4}.
\]
Equality is attained by \(u=(1,0)\) and \(v=(0,1)\), since both
\(\norm{u+\eps v}_4\) and \(\norm{u-\eps v}_4\) equal
\((1+\eps^4)^{1/4}\).
\end{proof}

It remains to exhibit one quasi-orthogonal pair with a smaller one-sign value
at \(\eps=1/2\).  Put
\[
s=\frac16,\qquad a=(1-s^4)^{1/4},\qquad x=(a,s).
\]
Then \(\norm{x}_4=1\).  Let
\[
r=\left(\frac{s}{a}\right)^3,\qquad c=(1+r^4)^{-1/4},\qquad y=c(r,-1).
\]
Then \(\norm{y}_4=1\).  The norming functional at \(x\) is
\[
p=(a^3,s^3)\in (\ell_4^2)^*=\ell_{4/3}^2,
\]
because \(\norm{p}_{4/3}=1\) and \(p(x)=a^4+s^4=1\).  Moreover
\[
p(y)=c(a^3r-s^3)=0.
\]
Thus \(y\urcorner x\) in the convention of the source paper.

We now bound the value of \(\norm{x+\frac12 y}_4\).  Since
\[
a^4=1-\frac1{6^4}=\frac{1295}{1296}>\left(\frac56\right)^4,
\]
we have \(a>5/6\), and hence
\[
0<r=\left(\frac{s}{a}\right)^3<\left(\frac{1/6}{5/6}\right)^3=\frac1{125}.
\]
Also \(0<c<1\).  Therefore the first coordinate of
\(x+\frac12 y\) satisfies
\[
0<a+\frac{cr}{2}<1+\frac1{250}=\frac{251}{250},
\]
and the second coordinate satisfies
\[
\left|s-\frac{c}{2}\right|<\frac13.
\]
Consequently
\[
\norm{x+\tfrac12 y}_4^4
<\left(\frac{251}{250}\right)^4+\left(\frac13\right)^4
=\frac{325405456081}{316406250000}.
\]
The last rational number is strictly smaller than \(17/16\), since
\[
\frac{17}{16}-\frac{325405456081}{316406250000}
=\frac{673511534}{19775390625}>0.
\]
Thus
\[
\zeta^-_{\ell_4^2}(1/2)
\leq \norm{x+\tfrac12 y}_4
< \left(\frac{17}{16}\right)^{1/4}
=\beta^-_{\ell_4^2}(1/2)+1.
\]
Subtracting \(1\) gives
\[
\zeta^-_{\ell_4^2}(1/2)-1<\beta^-_{\ell_4^2}(1/2),
\]
as claimed.

\section*{Verification notes}

The companion script \texttt{code/check\_l4\_counterexample.py} verifies the
rational inequality above and prints high-precision numerical values.  The
actual explicit pair gives approximately
\[
\norm{x+\tfrac12 y}_4=1.0051760632186776,
\qquad
(17/16)^{1/4}=1.0152715924344654.
\]

\section*{Novelty and literature check}

The following bounded searches were performed before promotion: the run's
registry, solutions, attempts, and proof-gap indexes were searched for
\texttt{1609.01587}, the title \texttt{New Moduli for Banach Spaces}, Milman
moduli, \(\beta^\pm\), \(\zeta^\pm\), quasi-orthogonality, Birkhoff-James
orthogonality, and related keywords.  Web searches were also run for the exact
paper title, \texttt{1609.01587}, \texttt{Milman's moduli}, \texttt{zeta beta
Banach}, and the quasi-orthogonality phrasing of the conjecture.  No later
answer to Conjecture 3 was found in this bounded search; the only direct hits
were the source paper and generic bibliographic entries.

\section*{Limitations}

This packet disproves only the equality
\[
\zeta^-_X(\eps)-1=\beta^-_X(\eps)
\]
from Conjecture 3, by the single example \(X=\ell_4^2\) and \(\eps=1/2\).  It
does not decide the proposed equality involving \(\zeta^+_X\) and
\(\beta^+_X\), nor the other questions and conjectures in the source paper.

\section*{Human-review recommendation}

This is a short, explicit counterexample with exact rational inequalities.
Human review should mainly confirm the source convention for
quasi-orthogonality and the elementary computation of \(\beta^-_{\ell_4^2}\).

\begin{thebibliography}{9}

\bibitem{IvanovMartini}
G. Ivanov and H. Martini,
\emph{New Moduli for Banach Spaces},
arXiv:1609.01587.

\bibitem{Milman}
V. D. Milman,
\emph{Geometric theory of Banach spaces. Part II. Geometry of the unit ball},
Russian Mathematical Surveys 26 (1971), 79--163.

\end{thebibliography}

\end{document}
